\unitopen{12}
\begin{candos}
\item opisać swój dzień po kolei (\de{zuerst, dann, danach, später});
\item użyć czasowników rozdzielnie złożonych w zdaniu i z czasownikiem modalnym;
\item wydać polecenie lub poprosić w trybie rozkazującym (\de{du}, \de{Sie}, \de{ihr});
\item przejść test końcowy A1: powitania, przedstawianie się, liczby, rodzajniki, biernik, czas, zakupy, modalne.
\end{candos}

\section{Słownictwo}
\vocabtable{12}

\section{Czasowniki rozdzielnie złożone}
Wiele czasowników ma przedrostek, który w zdaniu oznajmującym \emph{odpada i idzie na koniec}:
\de{aufstehen} $\rightarrow$ \de{Ich stehe um sieben Uhr \textbf{auf}.} \de{anfangen} $\rightarrow$ \de{Der Tag fängt um acht an.}
\de{einkaufen} $\rightarrow$ \de{Wir kaufen am Samstag \textbf{ein}.} W bezokoliczniku i z czasownikiem modalnym czasownik
jest w całości na końcu: \de{Ich muss früh \textbf{aufstehen}.} W pytaniu: \de{Rufst du mich \textbf{an}?}

\begin{center}\small
\begin{tabular}{@{}lll@{}}
\toprule
\textsf{\textbf{Bezokolicznik}} & \textsf{\textbf{Zdanie}} & \textsf{\textbf{Znaczenie}}\\
\midrule
\de{aufstehen} & \de{Ich stehe um 7 Uhr auf.} & wstaję\\
\de{ankommen} & \de{Der Zug kommt um 9 Uhr an.} & przyjeżdża\\
\de{anrufen} & \de{Ich rufe dich morgen an.} & zadzwonię\\
\de{mitkommen} & \de{Kommst du mit?} & idziesz z nami?\\
\de{einladen} & \de{Sie lädt uns ein.} & zaprasza nas\\
\bottomrule
\end{tabular}
\end{center}

\section{Tryb rozkazujący}
\begin{center}\small
\begin{tabular}{@{}lll@{}}
\toprule
\textsf{\textbf{Forma}} & \textsf{\textbf{Reguła}} & \textsf{\textbf{Przykład}}\\
\midrule
\de{du} & rdzeń, bez \de{-st} i bez \de{du} & \de{Komm!} \de{Mach das!} \de{Warte!}\\
\de{ihr} & forma \de{ihr}, bez \de{ihr} & \de{Kommt!} \de{Macht das!}\\
\de{Sie} & bezokolicznik $+$ \de{Sie} & \de{Kommen Sie!} \de{Machen Sie das!}\\
nieregularne & \de{e}$\rightarrow$\de{i}: bez końcówki & \de{Nimm!} \de{Sprich!} \de{Lies!} \de{Gib!}\\
\de{sein} & forma specjalna & \de{Sei ruhig!} \de{Seien Sie so freundlich!}\\
\bottomrule
\end{tabular}
\end{center}
Bez \de{bitte} rozkaz brzmi ostro: \de{Kommen Sie bitte herein.} W pracy częściej prosisz pytaniem (\de{Könnten Sie \dots?}).
Czasowniki rozdzielne: \de{Steh auf!} \de{Rufen Sie mich an!}

\begin{falle}
Polacy wstawiają przedrostek tuż po czasowniku: \de{Ich aufstehe um 7.} (błąd). Przedrostek zawsze
na końcu: \de{Ich stehe um 7 \textbf{auf}.} I druga rzecz: \de{anrufen} używa biernika: \de{Ich rufe \textbf{dich} an}
(\emph{dzwonię do ciebie}), nie \de{dir}. Wreszcie: \de{Ich warte \textbf{auf} den Bus} (czekam \emph{na} autobus: \de{auf} + biernik).
\end{falle}

\begin{szwajcaria}
Szwajcarzy kończą dzień pracy słowem \de{Feierabend}: \de{Schönen Feierabend!} mówi się koledze, gdy wychodzi
z biura. Dialektalne \de{Schaffe} znaczy „pracować”. W Szwajcarii ceni się \emph{pünktlichkeit} (pisz: \de{Pünktlichkeit})
i spotkania zaczynają się dokładnie o czasie.
\end{szwajcaria}

\begin{praca}
Typowy dzień programisty: \de{Um neun Uhr fängt das Daily an. Danach arbeite ich an einem Ticket. Mittags esse ich
mit den Kollegen. Am Nachmittag habe ich ein Review. Um halb sechs mache ich Feierabend.} Skopiuj ten szablon
i zmień go na swój dzień.
\end{praca}

\section{Dialogi}
\textbf{Dialog 1 --- co robisz w ciągu dnia?}
\begin{dialog}
\turn{A}{Wann stehst du normalerweise auf?}{O której zwykle wstajesz?}
\turn{B}{Ich stehe um halb sieben auf. Zuerst dusche ich und frühstücke.}{Wstaję o wpół do siódmej. Najpierw biorę prysznic i jem śniadanie.}
\turn{A}{Und wann fängt deine Arbeit an?}{A kiedy zaczyna się twoja praca?}
\turn{B}{Um neun. Ich komme um Viertel vor neun im Büro an.}{O dziewiątej. Przychodzę do biura za kwadrans dziewiąta.}
\turn{A}{Kommst du heute Abend mit ins Restaurant?}{Idziesz dziś wieczorem z nami do restauracji?}
\turn{B}{Gern! Ruf mich um sechs an.}{Chętnie! Zadzwoń do mnie o szóstej.}
\end{dialog}

\section{Ćwiczenia i test końcowy A1}
Wynik poniżej 80\% (11 z 14) oznacza powtórkę tygodni 5, 8 i 11 przed przejściem do A2.

\begin{cwiczenia}
\exitem{Połącz zdanie: \de{aufstehen} (ich), o siódmej.}{\de{Ich stehe um sieben Uhr auf.}}
\exitem{Połącz zdanie: \de{ankommen} (der Zug), um 9 Uhr.}{\de{Der Zug kommt um neun Uhr an.}}
\exitem{Ułóż zdanie z modalnym: \de{ich / muss / früh / aufstehen}}{\de{Ich muss früh aufstehen.}}
\exitem{Imperativ \de{du}: \de{kommen}.}{\de{Komm!}}
\exitem{Imperativ \de{Sie}: \de{warten}.}{\de{Warten Sie!}}
\exitem{Imperativ \de{du}: \de{nehmen}.}{\de{Nimm!}}
\exitem{Przedstaw się: imię, kraj, zawód (trzy zdania).}{np.\ \de{Ich heiße Piotr. Ich komme aus Polen. Ich bin Entwickler.}}
\exitem{Uzupełnij: \de{Ich habe \dots{} Bruder.} (ein, biernik)}{\de{einen}}
\exitem{Zaprzecz: \de{Ich habe Zeit.}}{\de{Ich habe keine Zeit.}}
\exitem{Podaj godzinę potocznie: 08:30.}{\de{halb neun}}
\exitem{Zamów kawę grzecznie.}{\de{Ich hätte gern einen Kaffee.} lub \de{Ich möchte einen Kaffee.}}
\exitem{Przetłumacz: „Czy możesz mi pomóc?”}{\de{Kannst du mir helfen?}}
\exitem{Podaj \de{Wo?}: \de{Der Laptop liegt \dots{} Tisch.}}{\de{auf dem}}
\exitem{Popraw: \de{Ich aufstehe um 7.}}{\de{Ich stehe um 7 auf.}}
\end{cwiczenia}

\section{Trening słownictwa (generowany)}
\drill{12}

\weekend{\textbf{Zadanie końcowe A1.} Opowiedz w 90 sekundach: kim jesteś, skąd pochodzisz, co robisz zawodowo, jak wygląda Twój
dzień i co chcesz robić w Szwajcarii. Nagraj. Odsłuchaj po tygodniu buforowym i porównaj z nagraniem z tygodnia~1.}

\begin{uwaga}[Tydzień buforowy A1]
Po tej jednostce jest tydzień buforowy: powtórz słownictwo z Anki (wszystkie karty A1), przepisz błędne zdania
z Dodatku~A, zrób dwa nagrania po 2 minuty, wyślij je do korekty. Dopiero potem zaczynasz A2.
\end{uwaga}
